\documentclass[9pt,a4paper]{article}

\usepackage[left=1.25cm,right=1.25cm,top=1.2cm,bottom=1.2cm]{geometry}
\usepackage[T2A]{fontenc}
\usepackage[utf8]{inputenc}
\usepackage[russian]{babel}

\usepackage{xcolor}
\usepackage[hidelinks]{hyperref}
\usepackage{enumitem}
\usepackage{tcolorbox}
\usepackage{fontawesome5}
\usepackage{titlesec}
\usepackage{parskip}
\usepackage{paracol}
\usepackage{ragged2e}

\tcbuselibrary{breakable}

\pagenumbering{gobble}
\setlength{\parindent}{0pt}
\setlist[itemize]{leftmargin=*, itemsep=1pt, topsep=1pt}

\definecolor{primary}{HTML}{0F172A}
\definecolor{accent}{HTML}{1E3A8A}
\definecolor{panel}{HTML}{F1F5F9}
\definecolor{text}{HTML}{0B1220}

\color{text}

\titleformat{\section}{\normalsize\bfseries\color{accent}}{}{0em}{}[\vspace{2pt}\hrule\vspace{6pt}]
\titlespacing*{\section}{0pt}{8pt}{6pt}

\begin{document}

%======================
% HEADER
%======================
\begin{tcolorbox}[
  colback=primary,
  colframe=primary,
  boxrule=0pt,
  arc=8pt,
  left=10pt,right=10pt,top=10pt,bottom=10pt
]
{\color{white}\Huge\textbf{Федосов Даниил Станиславович}}\\[-2pt]
{\color{white}\large System / Technical Analyst}\\[6pt]
{\color{white}
Москва \quad
\faPhone\ 8-915-513-38-85 \quad
\faEnvelope\ \href{mailto:fedosov.danil46@gmail.com}{fedosov.danil46@gmail.com} \quad
Telegram: @fedosoff46
}
\end{tcolorbox}

\vspace{8pt}

%======================
% BODY
%======================
\columnratio{0.33,0.64}
\begin{paracol}{2}
\RaggedRight

% LEFT COLUMN
\begin{tcolorbox}[
  breakable,
  colback=panel,
  colframe=panel,
  boxrule=0pt,
  arc=10pt,
  left=10pt,right=10pt,top=10pt,bottom=10pt
]

\section*{PROFILE | ПРОФИЛЬ}
\section*{PROFILE | ПРОФИЛЬ}
Системный / технический аналитик с опытом работы в enterprise и государственных ИС. 
Экспертиза: анализ инцидентов, диагностика систем через данные и логи, формализация требований 
и описание процессов взаимодействия систем.

\section*{KEY SKILLS | КЛЮЧЕВЫЕ НАВЫКИ}

\textbf{Системный анализ}
\begin{itemize}
  \item Сбор и формализация требований
  \item Декомпозиция задач
  \item Постановка задач в Jira
  \item Acceptance criteria / Definition of Done
  \item Анализ бизнес- и системной логики
\end{itemize}

\textbf{Данные и БД}
\begin{itemize}
  \item Реляционные БД: PostgreSQL, MySQL, Oracle
  \item NoSQL: MongoDB
  \item SQL: JOIN, CTE, агрегации, оконные функции
  \item Аналитические выборки и исследование данных
  \item Анализ жизненного цикла объектов и транзакций
  \item Диагностика системных ошибок через анализ данных
  \item Работа с лог-данными (ElasticSearch)
\end{itemize}

\textbf{Интеграции и API}
\begin{itemize}
  \item REST API, JSON
  \item OpenAPI / Swagger
  \item Postman
  \item Базовое понимание клиент-серверного взаимодействия
  \item Анализ контрактов и параметров запросов
\end{itemize}

\textbf{Процессы и моделирование}
\begin{itemize}
  \item BPMN 2.0
  \item Camunda
  \item Описание сценариев и точек отказа
  \item Анализ статусов, переменных и маршрутов процесса
\end{itemize}

\textbf{Диагностика и сопровождение}
\begin{itemize}
  \item ElasticSearch, Kibana
  \item Анализ логов и RCA
  \item Инциденты 2--3 линии
  \item SLA / приоритизация / эскалации
\end{itemize}

\textbf{Инструменты}
\begin{itemize}
  \item Jira, Confluence
  \item Git
  \item Python (скрипты и автоматизация)
  \item Scrum / Kanban
\end{itemize}

\section*{EDUCATION | ОБРАЗОВАНИЕ}
\textbf{Курский государственный университет}\\
Магистратура: 09.04.01 --- Информатика и вычислительная техника \hfill 2025 -- н.в.\\
Бакалавриат: 02.03.03 --- Математическое обеспечение и администрирование информационных систем \hfill 2021 -- 2025

\end{tcolorbox}

\switchcolumn

% RIGHT COLUMN
\section*{EXPERIENCE | ОПЫТ}

\textbf{ООО «ЭЙТИ КОНСАЛТИНГ»} \hfill 07.2025 --- н.в.\\
\textit{System Analyst (федеральная ГИС)}\\
\textit{Проект: государственная информационная система для обработки и мониторинга заявок и процессов ведомственного взаимодействия}
\begin{itemize}
  \item Анализ инцидентов 2--3 линии: воспроизведение сценариев, локализация причины, сбор технической фактуры и сопровождение до устранения
\item Анализ бизнес- и системной логики процессов в Camunda: маршруты, переменные, последовательности шагов, точки отказа
\item Анализ логов в ElasticSearch и Kibana, выявление повторяющихся паттернов ошибок
\item Анализ данных в PostgreSQL: SQL-запросы для диагностики, проверки гипотез и анализа состояния данных в продуктивном контуре
\item Формализация постановок задач для разработки в Jira: описание проблемы, сценарии воспроизведения, критерии готовности, декомпозиция
\item Взаимодействие с разработчиками и смежными командами при анализе дефектов и согласовании решений
  \item Контроль прохождения инцидентов по SLA до полного закрытия
  \item Проектирование логики обработки инцидентов и маршрутизации заявок в подсистеме: описание сценариев, статусов, переходов и условий выполнения процессов
  \item Проектирование логики обработки данных и сценариев взаимодействия модулей подсистемы
  \item \textbf{Автоматизация инцидентов:}
внедрил классификатор тикетов (NLP), снизив поток задач команды линии поддержки на \textbf{40\%}.
\item \textbf{Аналитика данных системы:}
через SQL-исследование выявил скрытые дефекты обработки данных и инициировал их исправление.
\end{itemize}

\textbf{ОКУ «УКС Курской области»} \hfill 07.2024 --- 12.2025\\
\textit{Ведущий специалист (цифровая трансформация)}\\
\textit{Проект: работа с государственной системой управления проектами капитального строительства (ИСУП Минстроя РФ)}
\begin{itemize}
  \item Сопровождение работы с ИСУП Минстроя РФ по объектам капитального строительства
  \item Контроль полноты и корректности данных, статусов и сроков исполнения
  \item Анализ отклонения и несоответствия в данных, инициализация корректировок
  \item Координация взаимодействие между муниципалитетами, подрядчиками и ответственными участниками процесса
  \item Обеспечение качества и актуальности данных для управленческой отчетности и контроля исполнения
\end{itemize}

\textbf{ООО «Информационный центр»} \hfill 08.2022 --- 07.2024\\
\textit{Developer / Team Lead}\\
\textit{Проект: разработка системы управления проектами и документооборотом INOUT Project}
\begin{itemize}
  \item Разработка и сопровождение ИСУП «INOUT Project»
  \item Координация задач команды: декомпозиция, планирование, контроль исполнения и ведение бэклога
  \item Взаимодействие с заказчиком по уточнению задач, ожиданий и приоритетов
  \item Подготовка релизов и сопровождение изменений
  \item Подготовка технической документации по системе: описание логики модулей, сценариев обработки данных и требований к доработкам
\end{itemize}

\section*{PROJECTS \& ADDITIONAL | ПРОЕКТЫ И ДОПОЛНИТЕЛЬНО}
\begin{itemize}
  \item \textbf{РФРИТ:} участие в подготовке грантовой заявки, проект прошёл в финальный этап отбора
  \item \textbf{Samsung IT Academy:} выпускник трека Android-разработки, опыт проектной работы с клиент-серверной логикой
  \item \textbf{Samsung IT Academy Hack:} победитель хакатона, лидер команды; отвечал за продуктовую упаковку и техническую реализацию прототипа
  \item \textbf{AR-проект для Минцифры:} участие в прототипировании решения
  \item \textbf{Наставник «Я в деле»:} сопровождение команд и доведение проектов до финальной защиты
\end{itemize}

\section*{POSITIONING | ПОЗИЦИОНИРОВАНИЕ}
\textbf{Интересующая роль:} System Analyst / Technical Analyst / Analyst в продуктовой или интеграционной команде.\\
\textbf{Фокус:} процессы, данные, инциденты, диагностика, требования, взаимодействие с разработкой.

\end{paracol}

\end{document}